\documentclass{article}
\usepackage{listings}
\usepackage{xcolor}
\usepackage{geometry}
\usepackage{amsmath}
\geometry{a4paper, margin=1in}

\definecolor{codegreen}{rgb}{0,0.6,0}
\definecolor{codegray}{rgb}{0.5,0.5,0.5}
\definecolor{codepurple}{rgb}{0.58,0,0.82}
\definecolor{backcolour}{rgb}{0.95,0.95,0.92}

\lstdefinestyle{mystyle}{
    backgroundcolor=\color{backcolour},
    commentstyle=\color{codegreen},
    keywordstyle=\color{magenta},
    numberstyle=\tiny\color{codegray},
    stringstyle=\color{codepurple},
    basicstyle=\ttfamily\footnotesize,
    breakatwhitespace=false,
    breaklines=true,
    captionpos=b,
    keepspaces=true,
    numbers=left,
    numbersep=5pt,
    showspaces=false,
    showstringspaces=false,
    showtabs=false,
    tabsize=2
}

\lstset{style=mystyle}

\title{Compiler Design Lab Report 3\\YACC-Based Arithmetic Expression Syntax Analyzer}
\author{Abdullah Al Jubaer Gem\\ID: 210041226\\Section: 2B}
\date{\today}

\begin{document}

\maketitle

\section{Objective}
Design and implement a YACC-based syntax analyzer for arithmetic expressions. The parser recognizes and evaluates expressions supporting addition, subtraction, multiplication, division, exponentiation, and parentheses, while detecting syntax errors.

\section{Expression Grammar}
The grammar encodes operator precedence structurally — lower-precedence operators sit higher in the hierarchy. Exponentiation binds tighter than multiplication/division, which bind tighter than addition/subtraction. Parentheses and unary minus are handled at the \texttt{primary} level.

\begin{align*}
\text{expr}    &\rightarrow \text{expr} + \text{term} \mid \text{expr} - \text{term} \mid \text{term} \\
\text{term}    &\rightarrow \text{term} * \text{factor} \mid \text{term} / \text{factor} \mid \text{factor} \\
\text{factor}  &\rightarrow \text{factor} \mathbin{\hat{}} \text{primary} \mid \text{primary} \\
\text{primary} &\rightarrow \texttt{(}\;\text{expr}\;\texttt{)} \mid \texttt{-}\;\text{primary} \mid \text{NUMBER}
\end{align*}

\section{Implementation}

\subsection{FLEX Lexer}

\begin{lstlisting}[language=C, caption=FLEX Lexer]
%{
#include "y.tab.h"
#include <stdlib.h>
%}

%%

[0-9]+(\.[0-9]+)?   {
                        yylval.dval = atof(yytext);
                        return NUMBER;
                    }

[ \t]               ;
\n                  return '\n';
"+"                 return '+';
"-"                 return '-';
"*"                 return '*';
"/"                 return '/';
"^"                 return '^';
"("                 return '(';
")"                 return ')';
.                   return yytext[0];

%%

int yywrap() { return 1; }
\end{lstlisting}

\subsection{YACC Parser}

\begin{lstlisting}[language=C, caption=YACC Parser]
%{
#include <stdio.h>
#include <stdlib.h>
#include <math.h>

void yyerror(const char *s);
int yylex();
%}

%union { double dval; }

%token <dval> NUMBER
%type  <dval> expr term factor primary

%%

input: input line | ;

line:
      '\n'
    | expr '\n'  { printf("Result = %g\n", $1); }
    ;

expr:
      expr '+' term  { $$ = $1 + $3; }
    | expr '-' term  { $$ = $1 - $3; }
    | term           { $$ = $1; }
    ;

term:
      term '*' factor  { $$ = $1 * $3; }
    | term '/' factor  { if ($3 == 0) { yyerror("Division by zero"); $$ = 0; }
                         else $$ = $1 / $3; }
    | factor           { $$ = $1; }
    ;

factor:
      factor '^' primary  { $$ = pow($1, $3); }
    | primary             { $$ = $1; }
    ;

primary:
      '(' expr ')'  { $$ = $2; }
    | '-' primary   { $$ = -$2; }
    | NUMBER        { $$ = $1; }
    ;

%%

void yyerror(const char *s) { printf("Syntax Error\n"); }

int main()
{
    printf("Enter Expression:\n");
    yyparse();
    return 0;
}
\end{lstlisting}

\section{Compilation Steps}

\begin{lstlisting}[language=bash, caption=Build and Run]
yacc -d 210041226_2B_Lab03.y
flex 210041226_2B_Lab03.l
gcc lex.yy.c y.tab.c -o parser -lm
./parser < input.txt
\end{lstlisting}

\section{Sample Input and Output}

\subsection{Input (\texttt{input.txt})}
\begin{lstlisting}
5+3-2
(10-3)+2
3*4+2
10/2-1
2^3
2^3^2
3.5*2+1.5
-5+3
(2+3)*(4-1)
\end{lstlisting}

\subsection{Output (\texttt{output.txt})}
\begin{lstlisting}
Enter Expression:
Result = 6
Result = 9
Result = 14
Result = 4
Result = 8
Result = 64
Result = 8.5
Result = -2
Result = 15
\end{lstlisting}

\section{Conclusion}
The YACC-based parser successfully evaluates arithmetic expressions with correct operator precedence. FLEX handles tokenization and passes tokens to YACC, which applies the grammar rules to compute results. Floating-point support and unary minus were implemented as extensions to the base grammar.

\end{document}
